\documentclass[10pt,a4paper,final,oneside,notitlepage]{article}
\usepackage[utf8]{inputenc} % Acentos, etc.
\usepackage[T1]{fontenc}    % Silabeo correcto de palabras acentuadas
\usepackage[spanish,es-tabla]{babel} % Idioma español con "Tabla" por defecto

% Márgenes según pautas CoNaIISI:
% Área de impresión: 17,5 cm de ancho por 22,54 cm de alto.
% Columnas de 8,25 cm con 8 mm de separación (8,25 * 2 + 0,8 = 17,3 cm).
% Margen superior e inferior de 3,0 cm en todas las páginas.
% Márgenes laterales: 1,85 cm en cada lado (21,0 cm - 1,85 * 2 = 17,3 cm).
\usepackage[a4paper, left=1.85cm, right=1.85cm, top=3.0cm, bottom=3.0cm]{geometry}
\setlength{\columnsep}{8mm}

% Tipografías según CoNaIISI:
% Texto principal y matemáticas: Times
\usepackage{mathptmx}
% Epígrafes de figuras y tablas: Helvética (Sans-Serif) 9 pts negrita
\usepackage[scaled=0.92]{helvet}
% Código fuente: Courier
\usepackage{courier}

% Paquetes matemáticos y utilidades
\usepackage{amsmath}
\usepackage{amsfonts}
\usepackage{amssymb}
\usepackage{titlesec}
\usepackage{multicol}
\usepackage{cite}        % Citas compactas [1-3], [2, 5]
\usepackage{booktabs}    % Tablas profesionales
\usepackage{multirow}    % Celdas multifila
\usepackage{array}       % Columnas de ancho fijo
\usepackage{setspace}
\usepackage{enumitem}
\usepackage{xcolor}
\usepackage{graphicx}
\usepackage{tikz}

\usepackage[hidelinks]{hyperref} % Hipervínculos en color negro

% Formato de párrafos según CoNaIISI:
% Sangría de 5 mm, completamente justificado, sin espacio adicional entre párrafos.
\setlength{\parindent}{5mm}
\setlength{\parskip}{0pt}

% Epígrafes (Captions) según CoNaIISI:
% Helvética 9 pts negrita.
% Mayúscula solo inicial: "Figura 1. ..." y "Tabla 1. ..."
% Tabla: título centrado por sobre la misma.
% Figura: referencias por debajo de la misma.
\usepackage{caption}
\DeclareCaptionFont{helv9pt}{\fontsize{9pt}{11pt}\selectfont\sffamily\bfseries}
\captionsetup{
  font=helv9pt,
  labelfont=helv9pt,
  labelsep=period
}
\captionsetup[table]{
  position=above,
  skip=6pt,
  justification=centering
}
\captionsetup[figure]{
  position=below,
  skip=6pt,
  justification=centering
}

% Formato de títulos de secciones según CoNaIISI:
% Primer Nivel: Times 12 pts negrita, inicial mayúscula, numerado ("1. Introducción"), ant. 24 pts, post. 12 pts.
\titleformat{\section}
  {\fontsize{12pt}{14.4pt}\selectfont\bfseries}
  {\thesection.\ }
  {0pt}
  {}
\titleformat{name=\section,numberless}
  {\fontsize{12pt}{14.4pt}\selectfont\bfseries}
  {}
  {0pt}
  {}
\titlespacing*{\section}{0pt}{24pt}{12pt}

% Segundo Nivel: Times 11 pts negrita, alineado izquierda, sin sangría, ant. 6 pts, post. 6 pts.
\titleformat{\subsection}
  {\fontsize{11pt}{13.2pt}\selectfont\bfseries}
  {\thesubsection.\ }
  {0pt}
  {}
\titleformat{name=\subsection,numberless}
  {\fontsize{11pt}{13.2pt}\selectfont\bfseries}
  {}
  {0pt}
  {}
\titlespacing*{\subsection}{0pt}{6pt}{6pt}

% Tercer Nivel: Times 10 pts negrita, alineado izquierda, ant. 12 pts, post. 6 pts.
\titleformat{\subsubsection}
  {\fontsize{10pt}{12pt}\selectfont\bfseries}
  {\thesubsubsection.\ }
  {0pt}
  {}
\titleformat{name=\subsubsection,numberless}
  {\fontsize{10pt}{12pt}\selectfont\bfseries}
  {}
  {0pt}
  {}
\titlespacing*{\subsubsection}{0pt}{12pt}{6pt}

% Parámetros de flotantes
\renewcommand{\topfraction}{0.92}
\renewcommand{\bottomfraction}{0.6}
\renewcommand{\textfraction}{0.08}
\renewcommand{\floatpagefraction}{0.7}
\setcounter{topnumber}{3}
\setcounter{totalnumber}{4}

\pagestyle{empty} % Sin número de página

% Notas al pie: Times 9 pts, espaciado simple
\renewcommand{\footnotesize}{\fontsize{9pt}{11pt}\selectfont}

\makeatletter

% Control de versión anónima para revisión por pares (doble ciego):
% \anonimafalse  -> Modo evaluación ciega (oculta autores, filiaciones y repo identificable)
% \anonimafalse -> Modo versión final aceptada (con autores, filiaciones y contactos)
\newif\ifanonima
\anonimafalse

% Metadatos
\def\email#1{\gdef\@email{#1}}
\def\@email{}

\def\universidad#1{\gdef\@universidad{#1}}
\def\@universidad{}

% Título Principal y Autores según CoNaIISI:
% En la primera página debe empezar en el centímetro 3,49 desde el borde superior de la hoja.
% Centrado, Times 14 pts negrita. Mayúscula la primera letra de sustantivos, verbos, etc.
% Dejar dos líneas en blanco (~24 pt) después del título.
% Autores: Times 12 pts estilo normal (sin negrita). Filiaciones: 12 pts cursiva. Email incluido.
% Dejar dos líneas en blanco (~24 pt) antes del texto principal.
\newcommand{\titulo} {
  \begingroup
  \centering
  \vspace*{0.49cm} % 3.0 cm margen superior + 0.49 cm = 3.49 cm al borde superior
  {\fontsize{14pt}{17pt}\selectfont\bfseries\@title\par}
  
  \vspace{24pt}
  
  \ifanonima
    {\fontsize{12pt}{14.4pt}\selectfont\textit{[Artículo remitido para evaluación por pares externos y doble ciego]}\par}
    \vspace{6pt}
    {\fontsize{10pt}{12pt}\selectfont\textit{(Información de autores y filiación omitida según las pautas de evaluación CoNaIISI)}\par}
  \else
    {\fontsize{12pt}{14.4pt}\selectfont\normalfont\@author\par}
    \ifx\@universidad\@empty\else
      \vspace{4pt}
      {\fontsize{12pt}{14.4pt}\selectfont\itshape\@universidad\par}
    \fi
    \ifx\@email\@empty\else
      \vspace{4pt}
      {\fontsize{12pt}{14.4pt}\selectfont\normalfont\@email\par}
    \fi
  \fi
  
  \vspace{24pt}
  \endgroup
}
\makeatother

% Resumen según CoNaIISI:
% Ubicado en la parte superior de la columna izquierda.
% Título "Resumen" con Times 12 pts negrita, centrado en la columna.
% Texto con Times 10 pts cursiva, interlínea simple, justificado.
% Dejar dos líneas en blanco después del resumen (~20 pt).
\newcommand{\resumen}[1]{%
  \begin{center}%
    {\fontsize{12pt}{14.4pt}\selectfont\bfseries Resumen\par}%
  \end{center}%
  \vspace{4pt}%
  {\fontsize{10pt}{12pt}\selectfont\itshape #1\par}%
  \vspace{12pt}%
}

% Palabras clave
\newcommand{\palabrasclave}[1]{%
  \noindent{\fontsize{10pt}{12pt}\selectfont\textbf{Palabras clave---}#1\par}%
  \vspace{16pt}%
}

\newcommand{\miniseccion}[1]{%
  \vspace{6pt}%
  {\fontsize{10pt}{12pt}\selectfont\textbf{#1}\par}%
  \vspace{3pt}%
}

% Entorno de figuras en multicols si se requiere
\newenvironment{Figure}
  {\par\medskip\noindent\minipage{\linewidth}}
  {\endminipage\par\medskip}

% Bibliografía según CoNaIISI: Times 9 pts, espaciado simple, citas entre corchetes
\newcommand{\bibliofont}{\fontsize{9pt}{11pt}\selectfont}

\let\OLDthebibliography\thebibliography
\renewcommand\thebibliography[1]{
  \OLDthebibliography{#1}
  \bibliofont
  \setlength{\parskip}{0pt}
  \setlength{\itemsep}{1pt plus 0.5pt}
}